\section{Euclidean Setting and Support Coordinates}

This section defines the fixed-normal Euclidean family studied throughout
the paper.

\subsection{Ambient geometry}

Let
\[
\mathbb E^3
\]
denote three-dimensional Euclidean space with its standard inner product.

Let
\[
D\subset\mathbb E^3
\]
be a regular dodecahedron centered at the origin. Its twelve faces occur in
six opposite pairs. Throughout the paper, the twelve face-normal directions
of \(D\) remain fixed. Only the perpendicular distances of the corresponding
supporting planes from the origin are allowed to vary.

\subsection{Opposite-face normal pairs}

Choose one outward unit normal from each opposite-face pair:
\[
\mathcal N
=
\{\mathbf n_1,\ldots,\mathbf n_6\}.
\]

The complete set of outward normal directions is
\[
\{\pm\mathbf n_1,\ldots,\pm\mathbf n_6\}.
\]

The numbering is arbitrary. Every rotational symmetry of the regular
dodecahedron permutes these six antipodal normal pairs.

\subsection{Support realization}

Let
\[
\mathbf h=(h_1,\ldots,h_6)\in\mathbb R_{>0}^6.
\]

For each \(i\), the two supporting planes normal to
\(\pm\mathbf n_i\) are
\[
\mathbf n_i\cdot\mathbf x=\pm h_i.
\]

They bound the slab
\[
S_i(h_i)
=
\left\{
\mathbf x\in\mathbb E^3:
|\mathbf n_i\cdot\mathbf x|\le h_i
\right\}.
\]

\begin{definition}[Support realization]
For
\[
\mathbf h\in\mathbb R_{>0}^6,
\]
define
\[
P(\mathbf h)
=
\bigcap_{i=1}^{6}S_i(h_i)
=
\left\{
\mathbf x\in\mathbb E^3:
|\mathbf n_i\cdot\mathbf x|\le h_i
\text{ for all }i
\right\}.
\]
The vector \(\mathbf h\) is the \emph{support register}.
\end{definition}

The set \(P(\mathbf h)\) is convex as an intersection of closed
half-spaces. It is centrally symmetric because
\[
\mathbf x\in P(\mathbf h)
\quad\Longrightarrow\quad
-\mathbf x\in P(\mathbf h).
\]

The selected normal directions span \(\mathbb E^3\), so the intersection is
bounded for every positive support register. Thus \(P(\mathbf h)\) is a
centrally symmetric convex polytope.

\subsection{The regular state}

There exists \(h_0>0\) such that
\[
D=P(h_0\mathbf 1),
\qquad
\mathbf 1=(1,1,1,1,1,1).
\]

Hence the regular dodecahedron is represented by six equal support numbers.
The positive uniform ray is
\[
U_+
=
\{h\mathbf 1:h>0\}.
\]

Section~4 proves that this ray is precisely the regular Euclidean similarity
family in the fixed-normal realization.

\subsection{Support space and admissibility}

The support register lies in the six-dimensional vector space
\[
\mathbb R^6.
\]

The positive orthant
\[
\mathbb R_{>0}^6
\]
is the ambient support domain. Positivity keeps every opposite plane pair
at positive distance from the origin.

Positivity alone does not preserve dodecahedral incidence. Large support
variations may make a defining inequality redundant, contract an edge, or
alter the vertex structure. Section~4 therefore restricts the geometric
results to the connected support chamber containing \(h_0\mathbf 1\) and
preserving the dodecahedral combinatorial type.

\subsection{Completeness of the support register}

Within the active fixed-normal family, the support register uniquely
determines the realization.

\begin{proposition}[Uniqueness of the support register]
Let
\[
\mathbf h,\mathbf k\in\mathbb R_{>0}^6.
\]
Assume that all twelve defining planes support genuine facets of both
\(P(\mathbf h)\) and \(P(\mathbf k)\). If
\[
P(\mathbf h)=P(\mathbf k),
\]
then
\[
\mathbf h=\mathbf k.
\]
\end{proposition}

\begin{proof}
For each \(i\), the active facet with outward normal \(\mathbf n_i\) has
support value
\[
\max_{\mathbf x\in P(\mathbf h)}
\mathbf n_i\cdot\mathbf x
=
h_i.
\]

The same support value equals \(k_i\) for \(P(\mathbf k)\). Equal convex
polytopes have equal support values in every direction, so
\[
h_i=k_i
\]
for all \(i\).
\end{proof}

Thus the six coordinates are complete for the declared fixed-normal,
origin-centered, active-facet family.

\subsection{Symmetry action}

Let \(\Gamma\) be the rotational symmetry group of the regular
dodecahedron. Each \(g\in\Gamma\) induces a permutation
\[
\pi_g\in S_6
\]
of the opposite-face pairs.

Define
\[
\rho(g)(h_1,\ldots,h_6)
=
(h_{\pi_g^{-1}(1)},\ldots,h_{\pi_g^{-1}(6)}).
\]

Then
\[
\rho(g_1g_2)
=
\rho(g_1)\rho(g_2),
\]
so \(\rho\) is a linear permutation representation on \(\mathbb R^6\).

Each \(\rho(g)\) preserves the standard inner product, coordinate sums,
the positive orthant, and the uniform vector \(\mathbf 1\).

The realization is equivariant:
\[
P(\rho(g)\mathbf h)
=
gP(\mathbf h).
\]

Thus coordinate permutation in support space agrees with Euclidean rotation
of the realized polytope.

\subsection{Scope}

The family \(P(\mathbf h)\) varies only the six common distances of
opposite supporting plane pairs. It keeps fixed

\begin{itemize}
\item the twelve regular dodecahedral normal directions;
\item the common center at the origin;
\item equality of the two distances within each opposite pair; and
\item central symmetry.
\end{itemize}

It excludes varying normal directions, unequal opposite supports,
translations, independent face rotations, arbitrary vertex displacement,
and non-centrally symmetric realizations.

No mechanical or dynamical meaning is assigned to the support coordinates.
They record Euclidean plane positions only.

The next section derives the canonical decomposition of the support register
into its uniform and zero-mean components.
